%
% $Id: slides.tex 4228 2006-06-21 21:55:12Z jjamor $
%
%
% Compilar a .pdf con LaTeX (pdflatex)
% Es necesario instalar Beamer (paquete latex-beamer en Debian)
%

%
% Gr�ficos:
% Los gr�ficos pueden suministrarse en PNG, JPG, TIF, PDF, MPS
% Los EPS deben convertirse a PDF (usar epstopdf)
%

\documentclass{beamer}
\usetheme{Warsaw}
\usebackgroundtemplate{\includegraphics[width=\paperwidth]{format/libresoft-bg.png}}
\usepackage[spanish]{babel}
\usepackage[latin1]{inputenc}
\usepackage{graphics}
\usepackage{amssymb} % Simbolos matematicos
\usepackage{url}

%\definecolor{libresoftgreen}{RGB}{162,190,43}
%\definecolor{libresoftblue}{RGB}{0,98,143}

%\setbeamercolor{titlelike}{bg=libresoftgreen}

%% Metadatos del PDF.
\hypersetup{
  pdftitle={Libre Software Communities},
  pdfauthor={Daniel Izquierdo Cort�zar},
  pdfcreator={GSyC/Libresoft},
  pdfproducer=PDFLaTeX,
  pdfsubject={nn},
}
%%


\AtBeginSection[]
{
  \begin{frame}<presentation>
    \frametitle{Index}
    \tableofcontents[current]
  \end{frame}
}


\begin{document}

\title{Management of Libre Software Communities}
\subtitle{Master on Free Software}
\institute{dizquierdo@gsyc.escet.urjc.es\\
GSyC/Libresoft}
\author{Daniel Izquierdo Cort�zar}
\date{\today}

\frame{
\maketitle
\begin{center}
\includegraphics[width=6cm]{format/gsyc-urjc}
\end{center}
}


% Si el titulo o el autor se quieren acortar para los pies de p�gina
% se pueden redefinir aqu�:
%\title{Titulo corto}
%\author{Autores abreviado}


%% LICENCIA DE REDISTRIBUCION DE LAS TRANSPAS
\frame{
~
\vspace{4cm}

\begin{flushright}
{\tiny
(cc) 2011 Daniel Izquierdo Cort�zar. \\
Some rights reserved. This document is distributed under the Creative \\
            Commons Attribution-ShareAlike 2.5 licence, available in \\
            http://creativecommons.org/licenses/by-sa/2.5/

%  Este documento (o uno muy similar) está disponible en \\
%  \url{http://gsyc.escet.urjc.es/~jjamor/}
}
\end{flushright}
}
%%



\section{Technical Report}

\begin{frame}
\frametitle{Who?}
\begin{itemize}
 \item SME company wondering about freeing some software
\end{itemize}
\end{frame}



\begin{frame}
\frametitle{Why?}
\begin{itemize}
\item For sure, we'd like to create community around it and attract new people.
\item Bussiness model similar to MySQL is our main goal.
\end{itemize}
\end{frame}


\begin{frame}
\frametitle{What we need}
\begin{itemize}
\item Good reasoning about freeing the software
\item Our Competitors nowadays
\item Technical Infrastructure needed
\item Licenses, we have heard about some BSD or GPL, but not sure!
\item Good practices in the source code in FLOSS environment
\item Others...
\end{itemize}
\end{frame}


\begin{frame}
\frametitle{Different Contexts}
\begin{itemize}
\item Product similar to Open Office (problem: parts of the source code are owned by another company)
\item Product similar to Zope/Plone
\item Product similar to Kopete or Gaim
\item Excelent IDE for Python
\item Micro-blogging
\item Mobile browser for Internet
\end{itemize}
\end{frame}





\end{document}
